\documentclass[a4paper,10pt]{article}
\usepackage[francais]{babel}
\usepackage{url}

\input{letter-icube-uds.sty}

\selectlanguage{francais}
\begin{document}


\name{St\'ephane \textsc{Genaud}}
\fonction{Professeur des Universités}
\telephoneillkirch{+33 (0)3688 54542}
\email{genaud@unistra.fr}

\begin{letter}[Lettre de soutien au projet Appsdeck]{à}{Jury du concours ``Emergences''}

Madame, Monsieur,

Mon domaine de recherche en informatique, au sein du laboratoire \textsc{Icube},
porte  sur le  parallélisme et  l'informatique distribuée.  A l'intérieur  de ce
grand domaine,  les \textit{clouds} focalisent une  attention toute particulière
dans  notre équipe  en  raison de  l'impact incontournable  qu'ils  ont sur  nos
usages, et la façon dont ils révolutionnent l'économie du numérique.

Le cloud introduit une rupture profonde avec le mode traditionnel de gestion des
applications en ligne. Il offre aux  clients la possibilité de réserver et gérer
lui-même les  ressources (calcul ou  stockage) dont  il a besoin,  moyennant une
facturation proportionnelle à l'utilisation. Pour  une entreprise, le recours au
cloud facilite grandement l'estimation du budget consacré à l'infrastructure car
il n'y a pas  d'investissement (et donc pas de maintenance non  plus) et le coût
d'utilisation, proportionnel  au temps ou  à la  quantité utilisés, est  connu à
l'avance. C'est  un avantage perçu comme  décisif dans de nombreux  cas et amène
beaucoup d'organisations à externaliser leurs applications dans des clouds.

Il est possible pour les clients d'exploiter les infrastructures ``nues''
\footnote{Ce  qu'on  appelle  IaaS  (Infrastructure-as-a-Service)},  ou  de  les
exploiter à  travers des services  additionnels dont l'objectif est  d'offrir un
accès  mieux  adapté  au  type  de  besoin du  client,  lui  faisant  gagner  en
productivité.  Appsdeck fait  partie de ces solutions. Elle  automatise un grand
nombre  des tâches  nécessaires pour  déployer une  application internet  sur un
serveur. 

Cependant,  de nombreuses  difficultés  doivent être  résolues  pour assurer  un
service de qualité quand le nombre de clients croît. Si les solutions techniques
déjà  mises   en  {\oe}uvre  dans   le  prototype  actuellement  en   test  sont
convaincantes, il n'en  reste pas moins que des études  devront être menées pour
comprendre les mécanismes de montées en charge et d'équilibrage de la charge sur
l'infrastructure. Nous  avons travaillé  dans mon  équipe sur  des problématique
proches, et  nous sommes  convaincus que notre  expertise pourra  aider Appsdeck
dans son  développement. Nous partageons  déjà avec les membres  d'Appsdeck, une
expérience commune:
\begin{itemize}
\item  Léo  \textsc{Unbekandt},  dans  son cursus  d'élève-ingénieur  à  l'école
  d'ingénieur en informatique ENSIIE, a accompli  un stage dans notre équipe, en
  prenant une part très active au développement du logiciel de courtage de Cloud
  Schlouder (\url{http://schlouder.gforge.inria.fr/}).
\item Léo est parti en dernière  année de l'ENSIIE (2013-2014) à l'Université de
  Cranfield pour  y obtenir le  Master en Computer  Science. Dans ce  cursus, il
  doit rédiger une thèse  de Master sous la direction d'un  collègue avce qui je
  co-signe des articles de recherche.  Cette expérience lui apportera une vision
  claire  des solutions  de  l'état de  l'art et  lui  permettra d'ébaucher  ses
  propres idées.
\item Yann \textsc{Klis} intervient régulièrement auprès des élèves ingénieurs de
 l'ENSIIE sur le thème des applications webs, dont il est spécialiste.
\end{itemize}

Je suis  convaincu de l'ambition  du projet  Appsdeck car la  solution technique
visée  (et démontrée  sur les  prototypes) par  ses porteurs  est d'une  qualité
évidente.   L'industrialisation  du  projet nécessitera  très  probablement  des
efforts   supplémentaires   en  ingénierie   et   également   en  recherche   et
développement. De par  la proximité de notre thème de  recherche avec le projet,
je suis  persuadé que nous pouvons  constuire une collaboration qui  n'aura rien
d'artificiel.  Pour ces  raisons, j'appporte mon soutien le plus  fort au projet
Appsdeck.



\end{letter}
\begin{flushright}
\includegraphics[width=.30\textwidth]{signgenaud.jpg}
\end{flushright}
%\newpage
%\bibliographystyle{plain}
%\bibliography{julien_blaise_ED}


\end{document}
